% introduction.tex -- revised introduction
\section{Introduction}
\label{sec:intro}

Re-implementing a JCL job stream on another platform is risky in a specific
way: the new job usually runs, but it behaves differently in situations the
migration did not consider. A step that used to be bypassed after a failure now
runs, a consumer reads a dataset before its producer has rewritten it, or a
renamed output is no longer found downstream. These rules are implicit in the
semantics of JCL~\cite{ibm_jcl_reference}. Structural dependency tools record
which step uses which dataset, but not the obligations (``step~B must not run if
step~A ends with RC\,$>$\,4'') that the new implementation must preserve.

RELIC extracts such obligations from JCL as explicit behavioural contracts
(in the sense of design by contract~\cite{meyer1992applying}) and reports the
contracts that a changed version of a job breaks or weakens. This paper reports
a development-stage evaluation of RELIC and of the evaluation instrument built
around it. It does \emph{not} report an independent validation.

\paragraph{Research questions.}
\begin{description}
  \item[RQ1 (recovery)] How well do the contracts RELIC extracts match facts
    derived from an execution-semantics simulator of JCL?
  \item[RQ2 (detection)] Given an original and a mutated job, how well does a
    contract-level diff separate behaviour-changing from behaviour-preserving
    mutants, compared with structural, text-based and language-model baselines?
  \item[RQ3 (benchmark validity)] Does the benchmark discriminate a
    contract-aware detector from trivial change detectors?
\end{description}

\paragraph{Contributions.}
\begin{enumerate}
  \item A mutation benchmark for JCL change checking with complete accounting of
    every candidate mutant (656 candidates; 566 evaluable; 90 excluded with
    stated reasons) and ground-truth labels from an execution-semantics
    simulator that is code-independent of the detector (but not
    author-independent).
  \item Diagnostic and negative findings: on this benchmark a text diff that
    ignores JOB-card parameters scores F1~0.9978, higher than every RELIC variant
    (0.9048, 0.9510, 0.9800); the evaluable negatives are almost entirely
    cosmetic; and the committed language-model comparison is selection-biased
    and not interpretable.
  \item Behavioural-contract extraction and differencing for JCL
    (TEMPORAL, RESOURCE, CONTROL\_FLOW) with documented scope limits and an
    ablation of four post-hoc fixes.
  \item A catalogue of threats to validity and a list of experiments that would
    be needed before any claim of effectiveness.
\end{enumerate}

\paragraph{What we do not claim.} We do not claim that RELIC outperforms a
normalised text diff or a language model, that the mutants represent real
migration defects, that the oracle matches z/OS behaviour, that the V2 figures
generalise beyond the three corpora, or that RELIC is compatible with
production JCL. Each of these is \textbf{[PENDING EXPERIMENT]} where it would
need new evidence (Section~\ref{sec:threats}).

\paragraph{Organisation.} Section~\ref{sec:related} places the work in the
literature; Section~\ref{sec:method} describes the benchmark and oracle;
Section~\ref{sec:arch} the tool; Section~\ref{sec:results} the results;
Sections~\ref{sec:discussion} and~\ref{sec:threats} the interpretation and
threats; Section~\ref{sec:conclusion} concludes.
